% problem-000335 7 碗形多环芳烃合成
\section*{7. 碗形多环芳烃合成}
（图：）

\subsection*{7-1}
⼀种合成⼼环烯的路线如下，请指出其中a\~h代表的反应条件和/或化合物。

（图：）

\subsection*{7-2}
写出原料A和第⼀个中间产物B的系统名称。

\subsection*{7-3}
由C到D的反应中，除⽬标分⼦外，还⽣成2种⼩分⼦副产物，写出它们的结构简式或名称。

\subsection*{7-4}
指出最终产物⼼环烯(R=H)及其⼆甲基化物 $\mathrm { ( R = C H _ { 3 } ) }$ 分⼦中的对称元素。

\subsection*{7-5}
⼀种合成花烯的路线如下，请指出其中i和j代表的主要试剂或条件。

（图：）

\subsection*{7-6}
化合物E的⽴体异构体中，只有⼀种异构体能经过烯交换反应得到F，请指出是哪⼀种？简述理由。

\subsection*{7-7}
实验中，H转变为I的反应只得到题给的⼀种⽴体异构体，简述理由。

\subsection*{7-8}
花烯G分⼦的碗-碗翻转能垒低：

（图：）

但I却难以翻转，简述理由。

\subsection*{7-9}
最近⽂献报道了I的10,15,20-三甲基取代衍⽣物。请画出该衍⽣物的结构，并指出有⽆⼿性？简述理由。
